% ===== §1 Introduction =====
\section{Introduction}
\label{sec:intro}

\noindent
This paper is the tenth in a series on the philosophy of intimacy and the theory of justice. The prior papers established a framework in which the goodness of a relation is read in the geometry of its cycles of value, the good cycle distinguished from the vicious by the sign of an accumulated phase, and they closed on a question they did not resolve: that the joint creation of value presupposes that two persons can jointly recognise what value is being created, and that the language by which a bond subsists intervenes in that recognition. The present paper addresses the question prior to the goodness of a cycle, namely how relation and value are constituted at all, and it does so by treating three matters together: the constitution of value, the mechanism by which two persons jointly apprehend it, and the role of symbolisation, of the making of language, in that apprehension.

The central thesis is stated here in advance of its development. Value is not a property antecedently resident in an object or a relation, and it is not the unsymbolisable remainder of the Real that any naming would falsify. Value is a phenomenon emergent on the coupled dynamics of the three Lacanian registers, the Real, the Symbolic, and the Imaginary, transversal across the three and localisable in none. Joint attention is the mechanism of its participatory generation: two persons do not detect an antecedent value but generate it in the manner of their attending. And symbolisation is internal to its sustenance rather than opposed to it: a symbolisation that sustains the coupling of the Symbolic with the Real, the poetic, regenerates the value, while a symbolisation that severs that coupling, the settling, extinguishes it. The making of language is, on this account, not the enemy of the nameless but the medium of the value's recurrent regeneration.

The method of the paper proceeds in order of grounding, and the order is substantive rather than expository. The paper begins by establishing the structure within which value is located, the three registers as coupled systems (\xref{sec:registers}); it describes, phenomenologically and without explanatory supplement, the datum of co-apprehension (\xref{sec:phenomenology}, \xref{sec:structure}); it derives the ontology that the datum requires (\xref{sec:ontology}); and only then does it admit the formal apparatus, the quantum-structural articulation of participatory measurement (\xref{sec:grammar}). It next confronts the limit of formal modelling directly, in a dilemma that conditions the entire formal programme (\xref{sec:modelling}); develops the political economy of the reflexively reconfigured subject (\xref{sec:political}); specifies the relation among the disciplinary frameworks as symmetry-broken phases of one structure (\xref{sec:breaking}); and concludes in a practice (\xref{sec:praxis}) and an ethics (\xref{sec:keystone}), terminating in plural conclusions rather than a synthesis (\xref{sec:conclusions}). The motivation for admitting the formal apparatus last is that each formal system is a settled symbolic apparatus imported from a domain in which its results are already cashed; to introduce it before the phenomenon is described is to expend, in advance, a determination the description has not underwritten, which is the operation the series identifies as settling failure. The apparatus is to be earned from the phenomenon and not expended upon it.

The contribution of the paper is fourfold. First, it supplies an ontology of value as an emergent of coupled registers, which yields a dynamical account of why the settling of value into pure symbol extinguishes rather than merely impoverishes it. Second, it articulates the participatory generation of value in the formal structure of quantum measurement, employed strictly as an epistemic model. Third, it isolates and examines a dilemma in the formal modelling of value: that the subject who creates value is reconfigured within the cycle of value, so that the generating mechanism is rewritten by its own products, with the consequence that the standard formal tools, the grammars of the Chomsky hierarchy and the fixed-law dynamical systems, are inadequate in a manner that is structural rather than incidental. Fourth, it grounds an ethics of the unguaranteeable good in an argument from the is--ought gap. The next section situates these contributions with respect to the frameworks on which they draw and against which they are defined.
